\oldpage[I]{84}
Supposons que $\varphi$ soit tel que tout $f\in A$ s'écrive
$f=h\varphi(f')$, où $h$ est inversible dans $A$ (ce qui est en particulier
le cas lorsque $\varphi$ est surjectif). Alors ${}^a\varphi$ est un
homéomorphisme de $X$ sur ${}^a\varphi(X)$.
\end{corollary}

Montrons que pour toute partie $E\subset A$, il existe une partie $E'$ de
$A'$ telle que $V(E)=V(\varphi(E'))$~; en vertu de l'axiome $(T_0)$
(\hyperref[I.1.1.8-fr]{1.1.8}) et de la formule
(\hyperref[I.1.2.2.1-fr]{1.2.2.1}), cela entraînera d'abord que
${}^a\varphi$ est injective, puis, toujours en vertu de
(\hyperref[I.1.2.2.1-fr]{1.2.2.1}), que ${}^a\varphi$ est un
homéomorphisme. Or, il suffit pour chaque $f\in E$ de prendre un $f'\in A'$ tel
que $h\varphi(f')=f$ avec $h$ inversible dans $A$~; l'ensemble $E'$ de ces
éléments $f'$ répond à la question.

\begin{env}[1.2.5]
\label{I.1.2.5-fr}
En particulier, lorsque $\varphi$ est l'homomorphisme canonique de $A$ sur un
anneau quotient $A/\mathfrak{a}$, on retrouve
(\hyperref[I.1.1.12-fr]{1.1.12}), et ${}^a\varphi$ est l'injection canonique
de $V(\mathfrak{a})$, identifié à $\operatorname{Spec}(A/\mathfrak{a})$,
dans $X=\operatorname{Spec}(A)$.
\end{env}

Un autre cas particulier de (\hyperref[I.1.2.4-fr]{1.2.4}) donne~:

\begin{corollary}[1.2.6]
\label{I.1.2.6-fr}
Si $S$ est une partie multiplicative de $A$, le spectre
$\operatorname{Spec}(S^{-1}A)$ s'identifie canoniquement (avec sa topologie)
au sous-espace de $X=\operatorname{Spec}(A)$ formé des $x$ tels que
$\mathfrak{j}_x\cap S=\varnothing$.
\end{corollary}

On sait en effet (\hyperref[0.1.2.6-fr]{0, 1.2.6}) que les idéaux premiers de
$S^{-1}A$ sont les idéaux $S^{-1}\mathfrak{j}_x$ tels que
$\mathfrak{j}_x\cap S=\varnothing$, et que l'on a
$\mathfrak{j}_x=({}^S i_A)^{-1}(S^{-1}\mathfrak{j}_x)$. Il suffit donc
d'appliquer à ${}^S i_A$ le cor. (\hyperref[I.1.2.4-fr]{1.2.4}).

\begin{corollary}[1.2.7]
\label{I.1.2.7-fr}
Pour que ${}^a\varphi(X)$ soit partout dense dans $X'$, il faut et il suffit
que tout élément du noyau $\operatorname{Ker}\varphi$ soit nilpotent.
\end{corollary}

En effet, appliquant la formule (\hyperref[I.1.2.2.3-fr]{1.2.2.3}) à l'idéal
$\mathfrak{a}=(0)$, il vient
$\overline{{}^a\varphi(X)}=V(\operatorname{Ker}\varphi)$, et pour que
$V(\operatorname{Ker}\varphi)=X$, il faut et il suffit que
$\operatorname{Ker}\varphi$ soit contenu dans tous les idéaux premiers de
$A$, c'est-à-dire dans le nilradical $\mathfrak{N}$ de $A$.

\subsection{Faisceau associé à un module}
\label{subsection:I.1.3-fr}
\label{I.1.3-fr}

\begin{env}[1.3.1]
\label{I.1.3.1-fr}
Soient $A$ un anneau commutatif, $M$ un $A$-module, $f$ un élément de $A$,
$S_f$ l'ensemble multiplicatif des $f^n$, où $n\geq0$. Rappelons que nous
posons $A_f=S_f^{-1}A$, $M_f=S_f^{-1}M$. Si $S'_f$ est la partie
multiplicative saturée de $A$ formée des $g\in A$ qui divisent un élément de
$S_f$, on sait que $A_f$ et $M_f$ s'identifient canoniquement à
${S'_f}^{-1}A$ et ${S'_f}^{-1}M$
(\hyperref[0.1.4.3-fr]{0, 1.4.3}).
\end{env}

\begin{lemma}[1.3.2]
\label{I.1.3.2-fr}
Les conditions suivantes sont équivalentes~:
\begin{enumerate}
  \item[\rm(a)] $g\in S'_f$~;
  \item[\rm(b)] $S'_g\subset S'_f$~;
  \item[\rm(c)] $f\in\mathfrak{r}(g)$~;
  \item[\rm(d)] $\mathfrak{r}(f)\subset\mathfrak{r}(g)$~;
  \item[\rm(e)] $V(g)\subset V(f)$~;
  \item[\rm(f)] $D(f)\subset D(g)$.
\end{enumerate}
\end{lemma}

Cela résulte immédiatement des définitions et de
(\hyperref[I.1.1.5-fr]{1.1.5}).

\begin{env}[1.3.3]
\label{I.1.3.3-fr}
Si $D(f)=D(g)$, le lemme (\hyperref[I.1.3.2-fr]{1.3.2, b)}) montre que
$M_f=M_g$. Plus généralement, si $D(f)\supset D(g)$, donc
$S'_f\subset S'_g$, on sait (\hyperref[0.1.4.1-fr]{0, 1.4.1}) qu'il existe un
homomorphisme canonique fonctoriel
\[
  \rho_{g,f}:M_f\longrightarrow M_g,
\]
et si $D(f)\supset D(g)\supset D(h)$, on a
(\hyperref[0.1.4.4-fr]{0, 1.4.4})
\begin{equation}
  \rho_{h,g}\circ\rho_{g,f}=\rho_{h,f}.
  \tag{1.3.3.1}\label{I.1.3.3.1-fr}
\end{equation}
\end{env}

